\section{Go-To eksempler til ligningsløsning med sin og cos}

Når vi skal løse trigonometriske ligninger, bruger vi ofte **sinus** og **cosinus**. Lad os tage nogle eksempler.

\subsection*{Eksempel 1: Beregning af en katete}
Vi ønsker at finde den hosliggende katete i en retvinklet trekant. Vi kender: 

\[
c = 2, \quad A = 30^\circ
\]

Her kan vi bruge **cosinusformlen**:

\[
b = c \cdot \Cos(A)
\]

Indsætter vi værdierne:

\[
b = 2 \cdot \Cos(30^\circ) \approx 1.732
\]

---

\subsection*{Eksempel 2: Beregning af vinkler}
Lad os nu sige, at vi kender begge kateter:

\[
a = 2, \quad b = 3, \quad c = \sqrt{a^2 + b^2} \approx 3.606
\]

Vi ønsker at finde vinklerne \(A\) og \(B\). Dette er en retvinklet trekant, så vi ved:

\[
C = 90^\circ
\]

Vi kan igen bruge cosinusformlen:

\[
b = c \cdot \Cos(A)
\]

For at isolere vinklen \(A\) omskriver vi ligningen:

\[
\frac{b}{c} = \Cos(A)
\]

Herefter bruger vi den **inverse cosinusfunktion**:

\[
A = \Cos^{-1}\left(\frac{b}{c}\right)
\]

Indsætter vi tallene:

\[
A = \Cos^{-1}\left(\frac{3}{3.606}\right) \approx 33.7^\circ
\]

Vinkel \(B\) findes ved at bruge, at summen af vinkler i en trekant altid er \(180^\circ\):

\[
B = 180^\circ - 90^\circ - 33.7^\circ \approx 56.3^\circ
\]

---

\textbf{Bemærkninger:}
\begin{enumerate}
    \item Vi bruger altid **inverse funktioner** for at isolere vinkler. 
    \item Husk at skrive grader (\(^\circ\)) når vi arbejder med vinkler.  
    \item Trin-for-trin-metoden: isoler trigonometrisk funktion → brug inverse funktion → indsæt tal.
\end{enumerate}
